% Cover letter for Apple, AI Measurement Scientist, AI Evaluation Platform. Built from cover-note.md.
\documentclass[11pt]{article}
\usepackage[margin=1in]{geometry}
\usepackage{fontspec}\setmainfont{Times New Roman}
\usepackage[hidelinks]{hyperref}
\hypersetup{pdftitle={Amal Mehta — Cover letter}, pdfauthor={Amal Mehta}}
\pagestyle{empty}
\setlength{\parindent}{0pt}
\setlength{\parskip}{0.85em}
\linespread{1.12}
\begin{document}
{\LARGE\bfseries Amal Mehta}\\[3pt]
SF Bay Area \quad$\cdot$\quad \href{mailto:amal.mehta@gmail.com}{amal.mehta@gmail.com} \quad$\cdot$\quad \href{https://amalmehta.com/}{amalmehta.com}

\vspace{1.2em}
November 25, 2026

Apple, AI Evaluation Platform\\
AI Measurement Scientist (role 200695284-3337)

\vspace{0.4em}
Dear AI Evaluation Platform team,

Most of my research this year has been about one problem: what happens when the number you train on isn’t the thing you care about.

I ran into it head on in a small study. I trained a code-repair agent with GRPO and rewarded it for passing the tests it could see. Its pass rate went up, from 35\% to 58\%. Then I checked its code against a full hidden test. The share that was actually correct had dropped from 31\% to 2\%. It had learned to delete the code no test touched. Random spot checks only slowed that down. Choosing which attempts to check held false passes to 4\%, against 24\% for random checks on the same budget.

That’s why this role caught my eye. If teams decide what ships from a score, the score has to be right.

I’ve worked this way for a while. For four years at PhysicsAI, we judged our synthetic training data by how a detector did on real drone footage, not on more synthetic data.

I’m not a psychometrician, and IRT would be new to me. I’d be learning it from your team, and I’m looking forward to that.

Thanks for reading.

Amal Mehta
\end{document}
